% Declare the main table early so the two-column float appears with the evaluation protocol.
\begin{table*}[!t]
 \centering\scriptsize\setlength{\tabcolsep}{3pt}\renewcommand{\arraystretch}{1.03}
 \begin{tabular}{@{}c@{\hspace{2pt}}|@{\hspace{2pt}}c@{\hspace{2pt}}|rr|rr|rr|rr|rr|rr@{}}
\toprule
\Lang{数据集}{Dataset} & \shortstack{\Lang{块缺失}{Block}\\\Lang{配置}{masking}} & \multicolumn{2}{c}{\shortstack{\Model{}\\\Lang{（本文）}{(Ours)}}} & \multicolumn{2}{c}{\shortstack{BRITS\\(2018)}} & \multicolumn{2}{c}{\shortstack{CSDI\\(2021)}} & \multicolumn{2}{c}{\shortstack{ImputeFormer\\(2024)}} & \multicolumn{2}{c}{\shortstack{TimesNet\\(2023)}} & \multicolumn{2}{c}{\shortstack{T1\\(2026)}} \\
\cmidrule(lr){3-4}\cmidrule(lr){5-6}\cmidrule(lr){7-8}\cmidrule(lr){9-10}\cmidrule(lr){11-12}\cmidrule(lr){13-14}
 & $p_s\,/\,b$ & MSE & MAE & MSE & MAE & MSE & MAE & MSE & MAE & MSE & MAE & MSE & MAE \\
\midrule
\multirow{4}{*}{\rotatebox[origin=c]{90}{ETTh1}} & $10\%\,/\,8$ & \textcolor{red}{0.1168} & \textcolor{red}{0.2055} & 0.5188 & 0.4990 & 0.1872 & 0.2473 & 0.2164 & 0.2914 & 0.3029 & 0.3517 & \textcolor{blue}{\underline{0.1323}} & \textcolor{blue}{\underline{0.2278}} \\
 & $15\%\,/\,8$ & \textcolor{red}{0.1722} & \textcolor{red}{0.2565} & 0.6529 & 0.5788 & 0.3038 & 0.3277 & 0.3057 & 0.3512 & 0.4208 & 0.4121 & \textcolor{blue}{\underline{0.1936}} & \textcolor{blue}{\underline{0.2805}} \\
 & $20\%\,/\,8$ & \textcolor{red}{0.2619} & \textcolor{red}{0.3227} & 0.7771 & 0.6397 & 0.3834 & 0.3945 & 0.4797 & 0.4410 & 0.5825 & 0.4817 & \textcolor{blue}{\underline{0.3125}} & \textcolor{blue}{\underline{0.3520}} \\
 & $10\%\,/\,16$ & \textcolor{red}{0.3350} & \textcolor{red}{0.3720} & 0.8855 & 0.6801 & 0.4517 & 0.4291 & 0.5322 & 0.5002 & 0.6397 & 0.5124 & \textcolor{blue}{\underline{0.3895}} & \textcolor{blue}{\underline{0.3987}} \\
\midrule
\multirow{4}{*}{\rotatebox[origin=c]{90}{Weather}} & $10\%\,/\,8$ & \textcolor{red}{0.0424} & \textcolor{blue}{\underline{0.0435}} & 0.0616 & 0.0685 & \textcolor{blue}{\underline{0.0427}} & \textcolor{red}{0.0361} & 0.0550 & 0.0547 & 0.0670 & 0.1062 & 0.0474 & 0.0646 \\
 & $15\%\,/\,8$ & \textcolor{blue}{\underline{0.0525}} & \textcolor{blue}{\underline{0.0539}} & 0.0719 & 0.0843 & \textcolor{red}{0.0506} & \textcolor{red}{0.0435} & 0.0636 & 0.0641 & 0.0829 & 0.1226 & 0.0608 & 0.0775 \\
 & $20\%\,/\,8$ & \textcolor{red}{0.0646} & \textcolor{blue}{\underline{0.0703}} & 0.0911 & 0.1079 & \textcolor{blue}{\underline{0.0665}} & \textcolor{red}{0.0589} & 0.0878 & 0.0808 & 0.1090 & 0.1466 & 0.0835 & 0.1041 \\
 & $10\%\,/\,16$ & \textcolor{blue}{\underline{0.0786}} & \textcolor{blue}{\underline{0.0929}} & 0.1057 & 0.1255 & \textcolor{red}{0.0776} & \textcolor{red}{0.0731} & 0.0919 & 0.0967 & 0.1621 & 0.1950 & 0.1237 & 0.1488 \\
\midrule
\multirow{4}{*}{\rotatebox[origin=c]{90}{Exchange}} & $10\%\,/\,8$ & \textcolor{blue}{\underline{0.0092}} & 0.0578 & 1.2295 & 0.8475 & 0.0098 & \textcolor{blue}{\underline{0.0575}} & 0.0576 & 0.1501 & 0.0196 & 0.0918 & \textcolor{red}{0.0084} & \textcolor{red}{0.0515} \\
 & $15\%\,/\,8$ & \textcolor{blue}{\underline{0.0154}} & \textcolor{blue}{\underline{0.0739}} & 1.2301 & 0.8567 & 0.0169 & 0.0752 & 0.0945 & 0.1931 & 0.0241 & 0.1006 & \textcolor{red}{0.0140} & \textcolor{red}{0.0662} \\
 & $20\%\,/\,8$ & 0.0542 & \textcolor{blue}{\underline{0.0985}} & 1.1269 & 0.8653 & \textcolor{red}{0.0447} & 0.1056 & 0.1195 & 0.2175 & 0.0577 & 0.1234 & \textcolor{blue}{\underline{0.0476}} & \textcolor{red}{0.0964} \\
 & $10\%\,/\,16$ & \textcolor{blue}{\underline{0.1623}} & \textcolor{red}{0.1720} & 1.2345 & 0.9078 & \textcolor{red}{0.1577} & 0.1864 & 0.2713 & 0.3147 & 0.1997 & 0.2001 & 0.1636 & \textcolor{blue}{\underline{0.1774}} \\
\midrule
\multirow{4}{*}{\rotatebox[origin=c]{90}{Beijing AQ}} & $10\%\,/\,8$ & \textcolor{red}{0.0772} & \textcolor{blue}{\underline{0.1446}} & 0.1023 & 0.1684 & \textcolor{blue}{\underline{0.0776}} & \textcolor{red}{0.1419} & 0.0868 & 0.1502 & 0.2273 & 0.2927 & 0.0861 & 0.1585 \\
 & $15\%\,/\,8$ & \textcolor{blue}{\underline{0.0948}} & \textcolor{blue}{\underline{0.1631}} & 0.1185 & 0.1843 & \textcolor{red}{0.0915} & \textcolor{red}{0.1575} & 0.0998 & 0.1658 & 0.3430 & 0.3638 & 0.1285 & 0.1971 \\
 & $20\%\,/\,8$ & \textcolor{blue}{\underline{0.1228}} & 0.1882 & 0.1426 & 0.2074 & \textcolor{red}{0.1176} & \textcolor{red}{0.1814} & 0.1244 & \textcolor{blue}{\underline{0.1877}} & 0.5200 & 0.4459 & 0.2255 & 0.2643 \\
 & $10\%\,/\,16$ & \textcolor{blue}{\underline{0.1468}} & \textcolor{blue}{\underline{0.2069}} & 0.1541 & 0.2150 & \textcolor{red}{0.1359} & \textcolor{red}{0.1950} & 0.1533 & 0.2101 & 0.7528 & 0.5501 & 0.3924 & 0.3617 \\
\midrule
\multirow{4}{*}{\rotatebox[origin=c]{90}{PhysioNet}} & $10\%\,/\,8$ & \textcolor{blue}{\underline{0.3868}} & \textcolor{blue}{\underline{0.3434}} & 0.5181 & 0.4190 & 0.5745 & 0.4547 & \textcolor{red}{0.3742} & \textcolor{red}{0.3337} & 0.5314 & 0.4412 & 0.4871 & 0.4144 \\
 & $15\%\,/\,8$ & \textcolor{blue}{\underline{0.4765}} & \textcolor{blue}{\underline{0.4038}} & 0.5882 & 0.4680 & 0.6721 & 0.5119 & \textcolor{red}{0.4614} & \textcolor{red}{0.3936} & 0.6574 & 0.5078 & 0.5970 & 0.4727 \\
 & $20\%\,/\,8$ & \textcolor{blue}{\underline{0.5648}} & \textcolor{blue}{\underline{0.4553}} & 0.6522 & 0.5090 & 1.0306 & 0.6609 & \textcolor{red}{0.5362} & \textcolor{red}{0.4440} & 0.7701 & 0.5579 & 0.7328 & 0.5394 \\
 & $10\%\,/\,16$ & \textcolor{blue}{\underline{0.5818}} & \textcolor{blue}{\underline{0.4699}} & 0.6863 & 0.5383 & 0.9781 & 0.6614 & \textcolor{red}{0.5494} & \textcolor{red}{0.4573} & 0.8291 & 0.5854 & 0.7736 & 0.5610 \\
\midrule
\multirow{4}{*}{\rotatebox[origin=c]{90}{PAMAP2}} & $10\%\,/\,8$ & \textcolor{red}{0.2276} & \textcolor{red}{0.1816} & 0.6593 & 0.4508 & 0.4340 & 0.2374 & 0.4385 & 0.2525 & 0.5484 & 0.3305 & \textcolor{blue}{\underline{0.2299}} & \textcolor{blue}{\underline{0.1867}} \\
 & $15\%\,/\,8$ & \textcolor{red}{0.3267} & \textcolor{red}{0.2308} & 0.7439 & 0.4847 & 0.5172 & 0.2796 & 0.5474 & 0.2979 & 0.6128 & 0.3548 & \textcolor{blue}{\underline{0.3367}} & \textcolor{blue}{\underline{0.2397}} \\
 & $20\%\,/\,8$ & \textcolor{red}{0.4669} & \textcolor{red}{0.2927} & 0.8260 & 0.5248 & 0.6143 & 0.3225 & 0.6146 & 0.3350 & 0.7205 & 0.3896 & \textcolor{blue}{\underline{0.5020}} & \textcolor{blue}{\underline{0.3066}} \\
 & $10\%\,/\,16$ & \textcolor{red}{0.5320} & \textcolor{red}{0.3346} & 0.9201 & 0.5719 & 0.6497 & \textcolor{blue}{\underline{0.3514}} & 0.6713 & 0.3693 & 0.7834 & 0.4232 & \textcolor{blue}{\underline{0.5913}} & 0.3536 \\
\bottomrule
\end{tabular}

 \caption{Imputation results across six datasets and four missingness configurations. Each training run averages twenty shared test masks; \Model{} reports the average over training seeds2023,2024 and 2025, while baselines retain their original training seeds. Best results are red; second-best results are blue and underlined.}
 \label{tab:formal}
\end{table*}

\section{Experiments}\label{sec:experiments}
\subsection{Datasets and Evaluation Protocol}
We evaluate \Model{} on ETTh1, Weather, Exchange-Rate, BeijingAirQuality, PhysioNet2012 and PAMAP2, covering industrial, meteorological, financial, environmental, clinical and wearable-sensor measurements. The prepared data use 7,21,8,12,35 and 40 variables, respectively. BeijingAirQuality treats the twelve stations as related channels. Native observation masks are retained; patient, subject and record boundaries are respected when constructing windows. All statistics used for scaling and history construction come from training data.

Artificial missingness is controlled by $p_s$ and $b$. The parameter $p_s$ controls the number of missing-event starts, while $b$ specifies how many consecutive time points each event hides. For each variable over a timeline of length $L$, we uniformly sample $\operatorname{round}(p_sL)$ distinct starts from positions that can accommodate a complete missing block. An event starting at $t$ hides the $b$ consecutive values from $t$ to $t+b-1$. Blocks from different starts may overlap, so the realized missing fraction depends jointly on $p_s$, $b$ and block overlap. For example, $p_s=10\%$ and $b=8$ sample a number of starts equal to $10\%$ of the timeline length and hide eight consecutive values per start. We evaluate four configurations: $10\%/8$, $15\%/8$, $20\%/8$ and $10\%/16$. Averaged across the six datasets and twenty test masks, these configurations produce approximately $57\%$, $73\%$, $83\%$ and $81\%$ artificial missingness, respectively, excluding native missingness.

For each dataset and missingness configuration, twenty fixed random seeds generate twenty distinct sets of missing positions using this procedure. All models are evaluated on the same twenty test masks for a fair comparison. Hash digests verify shared-mask identity, and the masks remain fixed across models and test outcomes. Native missingness is retained; MSE and MAE are computed only at originally observed values that are artificially hidden. Errors are measured in coordinates standardized using training-set statistics and averaged over the twenty test masks.

\subsection{Main Results}
Table~\ref{tab:formal} compares \Model{} with five baselines across six datasets and four masking configurations. Each model has adjacent MSE and MAE columns; rows group the four scenarios within each dataset. The long-block and higher-start-rate cases test the model with reduced visible support.
